\begin{abstract}
Marine biofouling is a coupled surface, biological, hydrodynamic, environmental, and maintenance problem. This critical review synthesizes an 81-paper curated corpus on antifouling and fouling-release coatings, nanomaterial-enabled surfaces, bioinspired and biobased interfaces, active cleaning, ship performance, and environmental risk. The evidence supports a distinction between preventing settlement, reducing adhesion, releasing established fouling, and physically removing fouling; these modes are not interchangeable and their performance depends on organism, exposure, motion, ageing, and cleaning conditions. Nanocomposites and functional interfaces frequently improve laboratory antifouling or adhesion metrics, but dispersion, wear, immersion stability, and constituent release remain deployment constraints. Ship-scale studies show that biofilm and roughness can impose drag and fuel penalties, while cleaning can introduce coating wear and contaminated effluent. The most defensible engineering pathway is therefore a system architecture combining low-adhesion surfaces, controlled active functions where justified, condition-based cleaning, hydrodynamic validation, and life-cycle assessment. The corpus remains heterogeneous, with limited long-duration full-scale and comparable ecotoxicological evidence.
\end{abstract}

\begin{IEEEkeywords}
marine biofouling, fouling-release coatings, nanocomposites, slippery surfaces, hull cleaning, hydrodynamic drag, environmental risk, life-cycle assessment
\end{IEEEkeywords}
